\section{The Resurrection of the Body}

To the ancient pagan world, no Christian teaching seemed stranger than the resurrection of the flesh. On the Areopagus in Athens, philosophers listened to St.\ Paul until he spoke of the resurrection of the dead, at which point some mocked, because Platonism sought the soul's escape from the prison of matter (Acts 17:32). Catholicism stands squarely against disembodied spiritualism. Because the eternal Word took true flesh from the Virgin Mary, died and rose in that flesh, and carried it to the Father's right hand, the redeemed are not ghosts permanently exiled from matter; heaven receives the whole human person, soul and risen body reunited forever.

Scripture lays this foundation across both Testaments. From his ash-heap Job cries out the prophecy the Church still chants beside Christian graves:

\begin{quote}
For I know that my Redeemer liveth, and in the last day I shall rise out of the earth. And I shall be clothed again with my skin, and in my flesh I shall see my God. Whom I myself shall see, and my eyes shall behold, and not another: this my hope is laid up in my bosom. (Job 19:25--27)
\end{quote}

In Ezechiel's valley of dry bones, sinews and flesh clothe the skeletons before the breath enters them: ``Behold I will open your graves, and will bring you out of your sepulchres, O my people'' (Ezech 37:1--14). Daniel foretells the universal awakening: ``And many of those that sleep in the dust of the earth, shall awake: some unto life everlasting, and others unto reproach, to see it always. But they that are learned, shall shine as the brightness of the firmament: and they that instruct many to justice, as stars for all eternity'' (Dan 12:2--3).

In 1 Corinthians 15, St.\ Paul binds our resurrection inseparably to Christ's: ``But now Christ is risen from the dead, the firstfruits of them that sleep: For by a man came death: and by a man the resurrection of the dead. And as in Adam all die, so also in Christ all shall be made alive'' (1 Cor 15:20--22). Like the bare grain of wheat dying in the furrow to rise as the full stalk (1 Cor 15:35--38), at the last trumpet ``this corruptible must put on incorruption: and this mortal must put on immortality. \dots\ Death is swallowed up in victory. O death, where is thy victory? O death, where is thy sting?'' (1 Cor 15:52--55). To the Philippians he names our exemplar: the Saviour ``will reform the body of our lowness, made like to the body of his glory, according to the operation whereby also he is able to subdue all things unto himself'' (Phil 3:21).

Against medieval dualists, the Fourth Lateran Council (1215) defined the numerical identity of the risen body: all men shall rise with their own bodies which they now bear (\textit{omnes cum suis propriis resurgent corporibus, quae nunc gestant}).\endnote{Fourth Lateran Council (1215), Constitution 1 (\textit{Firmiter credimus}), DS 801.} God does not hand the soul a substitute vessel; He raises the very flesh that knelt, fasted, laboured, received the Eucharist, and entered the grave.

The Fathers met physical objections with brisk realism. In his eighteenth Catechetical Lecture, St.\ Cyril of Jerusalem answers the mockers' puzzle of sailors eaten by fish, martyrs torn by beasts, or ashes scattered to the winds:

\begin{quote}
To thee, poor little feeble man, India is far from the land of the Goths, and Spain from Persia; but to God, who holds the whole earth in the hollow of His hand, all things are near at hand. \dots\ Imagine a mixture of seeds of different plants, and that these different seeds are contained in thy single hand; is it then to thee, who art a man, a difficult or an easy matter to separate what is in thine hand, and to collect each seed according to its nature, and restore it to its own kind? Canst thou then separate the things in thine hand, and cannot God separate the things contained in His hand, and restore them to their proper place?\endnote{St.\ Cyril of Jerusalem, \textit{Catechetical Lectures} 18.3, trans.\ Edwin Hamilton Gifford, NPNF2 vol.\ 7.}
\end{quote}

Every spring, bare vines and fig-trees quicken again as God's yearly sign against unbelief (\textit{Cat.}\ 18.7). Still more telling is moral justice: ``We do nothing without the body. We blaspheme with the mouth, and with the mouth we pray. \dots\ With the hand we rob, and by the hand we bestow alms. \dots\ Since then the body has been our minister in all things, it shall also share with us in the future the fruits of the past'' (\textit{Cat.}\ 18.19). Therefore, ``The root of all good works is the hope of the Resurrection,'' for ``he who believes that his body shall remain to rise again, is careful of his robe, and defiles it not with fornication'' (\textit{Cat.}\ 18.1).

Four centuries later, St.\ John of Damascus sums up Eastern patristic doctrine in Book IV, Chapter 27 of his \textit{Exposition of the Orthodox Faith}: ``By resurrection we mean resurrection of bodies. For resurrection is the second state of that which has fallen. For the souls are immortal, and hence how can they rise again? For if they define death as the separation of soul and body, resurrection surely is the re-union of soul and body, and the second state of the living creature that has suffered dissolution and downfall.''\endnote{St.\ John of Damascus, \textit{Exposition of the Orthodox Faith} IV.27, trans.\ S.\ D.\ F.\ Salmond, NPNF2 vol.\ 9 (Edinburgh: T.\ \& T.\ Clark, 1898).} Since the soul pursues neither virtue nor vice apart from the body, both share the reward, pledged already by Lazarus raised in his own flesh and by the Risen Lord offering His hands and side to Thomas.

In Book XXII of \textit{The City of God} (XXII.12--21), St.\ Augustine addresses infants, the aged, and the maimed: like a sculptor recasting a flawed statue from the exact same metal, God preserves every particle of natural substance while healing every deformity and blemish (\textit{City of God} XXII.19--20). Yet at the close of Chapter 19, Augustine pauses before one exception---the wounds of the martyrs:

\begin{quote}
But the love we bear to the blessed martyrs causes us, I know not how, to desire to see in the heavenly kingdom the marks of the wounds which they received for the name of Christ, and possibly we shall see them. For this will not be a deformity, but a mark of honour, and will add lustre to their appearance, and a spiritual, if not a bodily beauty. \dots\ While, therefore, it is quite true that no blemishes which the body has sustained shall appear in the resurrection, yet we are not to reckon or name these marks of virtue blemishes.\endnote{St.\ Augustine, \textit{The City of God} XXII.19, trans.\ Dods.}
\end{quote}

Augustine leaves the point open for the martyrs with modest reverence (``possibly we shall see them''), while illuminating why Christ Himself kept His five wounds on Easter Sunday. In \textit{ST} III, q.\ 54, a.\ 4, St.\ Thomas Aquinas cites Augustine and Bede on Luke 24:40 to give five reasons: as ``an everlasting trophy of His victory''; to confirm the disciples' faith; ``that when He pleads for us with the Father, He may always show the manner of death He endured for us''; to display to the redeemed the price of their mercy; and at the Judgment to convict the impenitent.\endnote{St.\ Thomas Aquinas, \textit{Summa Theologiae} III, q.\ 54, a.\ 4.} Those glorious scars add ``a special comeliness'' to the Risen Lord; what love endured for God is transfigured into splendor.

Yet if the purified soul already enjoys the Beatific Vision immediately after death (\textit{Benedictus Deus}), why does she still long for her body? In \textit{ST} I-II, q.\ 4, a.\ 5, Aquinas affirms that the separated soul truly enjoys essential beatitude before the Last Day, since beholding the Divine Essence requires no bodily phantasm (2 Cor 5:8). Yet the human soul is by nature the form of a body; apart from it, she lacks natural completeness. ``For the soul desires to enjoy God in such a way that the enjoyment also may overflow into the body, as far as possible,'' Aquinas explains. ``And therefore, as long as it enjoys God, without the fellowship of the body, its appetite is at rest in that which it has, in such a way, that it would still wish the body to attain to its share.''\endnote{St.\ Thomas Aquinas, \textit{Summa Theologiae} I-II, q.\ 4, a.\ 5, ad 4--5; cf.\ q.\ 4, a.\ 6.}

On one subtle point, Catholic theology notes a development within St.\ Thomas himself. In his early \textit{Sentences} commentary---preserved in the \textit{Supplement} (q.\ 93, a.\ 1)---he taught that reunion with the body increases beatitude in both \textit{extent} and \textit{intensity}.\endnote{\textit{Supplement}, q.\ 93, a.\ 1 (English Dominicans, 1922), where the editors explicitly note Aquinas's later retraction in \textit{ST} I-II, q.\ 4, a.\ 5, ad 5.} In the mature \textit{Summa Theologiae} (\textit{ST} I-II, q.\ 4, a.\ 5, ad 5), he expressly revised this judgment: because the separated soul is already wholly at rest as to the Object beheld, ``after the body has been resumed, Happiness increases not in intensity, but in extent,'' as ``the vigor of incorruption overflows into the lower nature'' (q.\ 4, a.\ 6).

St.\ Bernard of Clairvaux sets this truth to music in Chapter 11 of \textit{On the Love of God}. Even immersed ``in the fathomless depths of eternal light, and of a luminous eternity,'' the separated saints cherish a natural tenderness for the bodies that shared their pilgrimage:

\begin{quote}
The body then is for the soul a faithful companion; if it be a burden, it is also a help; when it ceases to aid, it also ceases to hinder; when again it helps, it is no more a burden. The first state is laborious, but profitable; the second inactive, but in no sense tedious; and the third is glorious. Hear how the Bridegroom in the Canticles invites the soul to these three conditions: ``Eat, O friends, and drink, and be inebriated, my dearly beloved.''\endnote{St.\ Bernard of Clairvaux, \textit{On the Love of God}, ch.\ 11, trans.\ Patmore (1884), pp.\ 47--52.}
\end{quote}

In this mortal life, Bernard explains on Cant 5:1, the soul eats the bread of daily obedience; after death, freed from corruptible flesh, she drinks long draughts of divine love mingled with quiet longing for her sleeping body; at the Resurrection, when her companion is restored in glory, the Bridegroom pours out the unmixed wine of the kingdom: ``In this life, eat; after death you shall drink; in the kingdom of the resurrection you shall be inebriated; you whom I well may call my dearly beloved.''
